% TOP_PROBLEM: {"id": 1239, "title": "Gauss Circle Problem", "category_id": 1, "category": {"id": 1, "name": "number_theory", "display_name": "Number Theory", "description": "Properties of integers, prime numbers, Diophantine equations.", "slug": "number-theory", "order_index": 1, "created_at": "2026-07-31T15:26:25.670Z"}, "status": "open"}
% FIRST_POSTED: null
% ENTRY_KIND: statement_only
% Imported from: batch05.tex; UnsolvedMath ID 1239
\documentclass[11pt]{article}
\usepackage[margin=25mm]{geometry}
\usepackage{fontspec}
\setmainfont{Latin Modern Roman}
\usepackage{amsmath,amssymb,mathtools}
\usepackage{unicode-math}
\setmathfont{Latin Modern Math}
\usepackage{xurl,hyperref}
\hypersetup{colorlinks=true,urlcolor=blue}
\setcounter{secnumdepth}{0}
\setlength{\parindent}{0pt}
\setlength{\parskip}{5pt}
\setlength{\emergencystretch}{3em}
\newcommand{\sref}[2]{\hyperlink{source:#1:#2}{[S#2]}}
\newcommand{\cataloguescope}[1]{\paragraph{Recorded target.}#1}
\begin{document}
% UnsolvedMath ID 1239; source catalogue code NT-032
\setcounter{section}{1238}
\section{Gauss Circle Problem}\label{1239}
{\small\sffamily UnsolvedMath ID 1239 · Number Theory}
\subsection{Definitions and mathematical statement}

\paragraph{Shared notation.}$\mathbb N=\{1,2,\ldots\}$, and $\mathbb Z,\mathbb Q,\mathbb R,\mathbb C$ denote the integers, rationals, reals and complex numbers. Any other conventions are specified locally.


For a real radius $r\geq1$, count lattice points on as well as inside the circle:
\[
N(r)=\#\{(a,b)\in\mathbb Z^2:a^2+b^2\leq r^2\},\qquad E(r)=N(r)-\pi r^2.
\]
For a real exponent $\alpha$, the estimate $E(r)=O_{\varepsilon}(r^{\alpha+\varepsilon})$ means that for every $\varepsilon>0$ there exist $C_\varepsilon>0$ and $r_\varepsilon\geq1$ with $|E(r)|\leq C_\varepsilon r^{\alpha+\varepsilon}$ for all $r\geq r_\varepsilon$.
\paragraph{Statement recorded in the supplied source.}
Determine the sharp growth of the circle-counting error $E(r)$. In particular, is
\[
\forall\varepsilon>0,\qquad E(r)=O_\varepsilon(r^{1/2+\varepsilon})\quad(r\to\infty)?
\] \sref{1239}{2}

\subsection{Short English statement}
How small can the error be when the area of a large circle approximates its number of lattice points? Does every exponent slightly greater than one-half give a valid error bound?
\subsection{Sources}
\begin{description}
\item[\hypertarget{source:1239:1}{S1}] ULAM.AI and the UnsolvedMath Contributors, UnsolvedMath: A Curated Collection of Open Mathematics Problems, record 1239 (NT-032). CC BY 4.0. \url{https://huggingface.co/datasets/ulamai/UnsolvedMath}.
\item[\hypertarget{source:1239:2}{S2}] Xiaochun Li and Xuerui Yang, \emph{An improvement on Gauss\textquotesingle s Circle Problem and Dirichlet\textquotesingle s Divisor Problem}, arXiv:2308.14859v2 (2023), Section 1, equation (1.1) and ensuing discussion. \url{https://arxiv.org/pdf/2308.14859}.
\end{description}
\label{end:1239}
\end{document}
